\documentclass[10pt,a4paper]{amsart}
\usepackage[margin=19mm]{geometry}
\usepackage{amsmath,amssymb,mathtools}
\usepackage[scr=rsfs,cal=euler]{mathalfa}
\usepackage{xfrac}
\usepackage[unicode,hidelinks,bookmarks=false]{hyperref}
\newcommand{\ie}{i.e.\ }
\newcommand{\cf}{cf.\ }
\newcommand{\resp}{respectively\ }
\newcommand{\Subsection}{\subsection}
\providecommand{\nobreakdash}{\nobreak\hbox{-}}
\setlength{\emergencystretch}{2em}
\multlinegap0pt
\usepackage{bm}
\usepackage{bbm}
\usepackage{dsfont}
\usepackage{xspace}
\usepackage{cancel}
\usepackage{mathtools}
\usepackage{amsthm}
\makeatletter
\@ifundefined{assumption}{\newtheorem{assumption}{Assumption}}{}
\@ifundefined{theorem}{\newtheorem{theorem}{Theorem}}{}
\@ifundefined{proposition}{\newtheorem{proposition}[theorem]{Proposition}}{}
\makeatother
\newcommand{\modifyfirst}[1]{#1}
\title[A separable and asymptotic-preserving dynamical low-rank met]{A separable and asymptotic-preserving dynamical low-rank method for the Vlasov-Poisson-Fokker-Planck system}
\author{Shiheng Zhang, Jingwei Hu}
\date{Source: arXiv:2601.11900v2 (CC BY 4.0)}
\begin{document}
\maketitle

\noindent\emph{Source and licence.} A separable and asymptotic-preserving dynamical low-rank method for the Vlasov-Poisson-Fokker-Planck system. Version arXiv:2601.11900v2 (CC BY 4.0), distributed under the Creative Commons Attribution 4.0 International licence, as recorded in the arXiv licence metadata for this exact version and shown on the abstract page https://arxiv.org/abs/2601.11900v2. Full rights notice: licenses/arxiv-2601.11900-CC-BY-4.0.txt.\par\smallskip

\noindent\textit{Editorial scope.} Counted unit: the Theorem whose source label is \texttt{thm:main\_AP} in the article (source0.tex lines 721--735, source label \texttt{thm:main\_AP}), delivered together with the article's own complete original proof of it (source0.tex lines 738--771, one \texttt{proof} environment of 34 lines). No mathematics is written, repaired or re-derived: statement and proof are the article's own text line for line. Required context retained uncounted: eq:L\_fd\_imex (lines 421--423); prop:coercivity (lines 605--621); prop:fluctuation\_bound (lines 662--673); asm:small\_fluct (lines 713--719). Every definition, display and label of those lines is retained unchanged. Omitted from this package and not redistributed: 1--420, 424--604, 622--661, 674--712, 720--720, 736--737, 772--1183. Nothing inside a retained excerpt is omitted or altered: every retained body is delivered exactly as the article's lines are written, so no omission receipt of any kind accompanies this package. Citations: the retained excerpts carry no citation command at all, so no bibliography entry of the article is transcribed and no reference is redistributed. Reference adaptation: none was needed and none was made; every reference inside the delivered unit resolves inside it, so no unresolved reference or citation is produced. Printed slips kept literal: no printed word, symbol, hypothesis or number of the counted statement or of its proof was altered. Layout and class: the article's own class is \texttt{article} with packages including graphicx, amsfonts, amsmath, amssymb, fancyhdr, titlesec, indentfirst, booktabs, verbatim, color, amsthm, hyperref, subcaption, appendix; the wrapper is the lane's pinned \texttt{amsart} editorial wrapper at 10pt on A4, which restates the theorem-like environments the retained text needs and adds the article's own macro definitions that the retained text needs (\texttt{\textbackslash modifyfirst}), with extra wrapper packages (bm, bbm, dsfont, xspace, cancel, mathtools). The delivered numbering is the unit's own. Assets: no retained span contains a figure environment, a graphics-inclusion command, a PDF-page inclusion, a table or a TikZ picture; no image file is copied into this package, the package declares no asset dependency and no image file is redistributed with it. Licence: version arXiv:2601.11900v2 of this article is distributed under the Creative Commons Attribution 4.0 International licence, as recorded in the arXiv licence metadata for this exact version and shown on the abstract page https://arxiv.org/abs/2601.11900v2. A full rights notice is delivered at licenses/arxiv-2601.11900-CC-BY-4.0.txt. Characters: every retained excerpt is pure ASCII, so the delivered engine, XeLaTeX with its default Latin Modern text font, reports no missing character.
\begin{equation}\label{eq:L_fd_imex}
  \frac{L_{i}^{n+1} - L_{i}^{(2)}}{\Delta t} = \langle \mathcal{A}_h(f^{(2)}), X_{i}^{n+1}\rangle_x^h + \frac{1}{\varepsilon}\langle \mathcal{L}_h^{n+1}(f^{n+1}), X_{i}^{n+1}\rangle_x^h,
\end{equation}

\begin{proposition}[Discrete coercivity]
\label{prop:coercivity}
Consider the decomposition $f = \rho M + g$ \modifyfirst{with $\rho=\langle f,1\rangle_v^h$, so that the fluctuation $g$ satisfies the zero velocity moment condition $\langle g, 1 \rangle_v^h = \mathbf{0}$}. There exists a constant $\gamma_h > 0$ such that the Fokker--Planck operator satisfies the coercivity estimate:
\begin{equation}
  -\big\langle \mathcal L_{h}(f),\, g/M \big\rangle_{xv}^h \;\ge\; \gamma_h\,\big\| g/M \big\|_{xv}^{h,2},
\end{equation}
where the constant $\gamma_h$ is defined as:
\begin{equation}
  \gamma_h:=\frac{M_{\min}\,\lambda_N}{\Big(1+\sqrt{{L_v\,M_{\max}}}\Big)^2}, \qquad \text{with} \quad \lambda_N=\frac{4}{\Delta v^2}\sin^2\!\Big(\frac{\pi}{2N_v}\Big),
\end{equation}
and the Maxwellian bounds
\[
  M_{\min}:=\min_{p,q} M_{p,q+1/2},
  \qquad
  M_{\max}:=\max\!\Big\{ \max_{p,q} M_{p,q},\, \max_{p,q} M_{p,q+1/2}\Big\}.
\]
\end{proposition}

\begin{proposition}[Fluctuation--based projection bound]
\label{prop:fluctuation_bound}
Let $P_X$ be the spatial projector onto the subspace $\mathrm{span}\{X_i\}$. Let $\bar{\beta}$ denote the arithmetic mean of $\beta_p$ over the spatial domain, and define the fluctuation $\delta \beta_p := \beta_p - \bar{\beta}$ and the inverse relative amplitude $C_\beta := \max_p (\beta_p \bar{\beta})^{-1}$.
Then, for any distribution $f$ in the range of $P_X$, the orthogonal projection of the Fokker--Planck operator is bounded by:
\begin{equation}
  \|(I-P_X)\mathcal L_{h}(f)\|_{xv}^h \le \kappa_h \|f\|_{xv}^h,
\end{equation}
where the fluctuation constant is defined by:
\begin{equation}\label{eq:kappa_h_def}
  \kappa_h \;:=\; \|\delta \beta\|_\infty \left( \frac{2\alpha_{\max}}{\Delta v^2} \;+\; C_\beta \frac{2}{\alpha_{\min}\Delta v^2} \right).
\end{equation}
\end{proposition}

\begin{assumption}[Small--fluctuation regime]\label{asm:small_fluct}
We assume that the field fluctuation $\kappa_h$ (defined in Proposition~\ref{prop:fluctuation_bound}) satisfies
\begin{equation}\label{eq:regime_condition}
    \kappa_h \, M_{\max}^{n+1} \;<\; \gamma_h,
\end{equation}
where $\gamma_h$ is the coercivity constant from Proposition~\ref{prop:coercivity} and $M_{\max}^{n+1} := \max\!\big\{ \max_{p,q} M_{p,q}^{n+1},\, \max_{p,q} M_{p,q+1/2}^{n+1}\big\}$, an upper bound that controls both the cell-centered and half-grid Maxwellian values used in the proof.
\end{assumption}

\begin{theorem}[AP property up to spatial fluctuation]
\label{thm:main_AP}
Let $f^{n+1}$ be the low--rank solution obtained after the L--step update. Let $\rho^{n+1} = \langle f^{n+1}, 1 \rangle_v^h$ be the macroscopic density vector. Let $G^{n+1}$ denote the discrete residual defined by the time difference and advection operators (see \eqref{eq:G_derivation}). Under Assumption~\ref{asm:small_fluct} and the formal asymptotic regime assumptions (implying $\|G^{n+1}\|_{xv}^h = \mathcal{O}(1)$), let $\Theta := \gamma_h / (\gamma_h - \kappa_h M_{\max}^{n+1}) \ge 1$. Then, the Fokker--Planck residual satisfies
\begin{equation}\label{eq:residual_bound}
  \big\|\mathcal L_{h}^{n+1}(f^{n+1})\big\|_{xv}^h
  \;\le\;
  \Theta \Big( \varepsilon\,\|G^{n+1}\|_{xv}^h \;+\; \kappa_h\,M_{\max}^{n+1}\,\sqrt{L_v}\,\big\|\rho^{n+1}\big\|_x^h \Big),
\end{equation}
and 
\begin{equation}\label{eq:distance_bound}
  \big\|f^{n+1}-\rho^{n+1} M^{n+1}\big\|_{xv}^h
  \;\le\;
  \frac{M_{\max}^{n+1}}{\gamma_h}\,\Theta \Big( \varepsilon\,\|G^{n+1}\|_{xv}^h \;+\; \kappa_h\,M_{\max}^{n+1}\,\sqrt{L_v}\,\big\|\rho^{n+1}\big\|_x^h \Big).
\end{equation}
\end{theorem}

\begin{proof}
We focus our analysis on the L-step because it determines the final state $f^{n+1}$. Recall the L-step update \eqref{eq:L_fd_imex} projected onto the updated spatial basis ${X}_i^{n+1}$:
\[
  \frac{L_i^{n+1} - L_i^{(2)}}{\Delta t} = \langle \mathcal{A}_h(f^{(2)}), {X}_i^{n+1} \rangle_x^h + \frac{1}{\varepsilon} \langle \mathcal{L}_{h}^{n+1}(f^{n+1}), {X}_i^{n+1} \rangle_x^h.
\]
Multiplying by ${X}_i^{n+1}$ and summing over $i$, we reconstruct the full distribution update projected onto the spatial subspace. Let $P_X^{n+1}$ be the orthogonal projector onto $\mathrm{span}\{{X}_i^{n+1}\}$. Rearranging the update equation to isolate the Fokker--Planck term yields:
\begin{equation}\label{eq:G_derivation}
  P_X^{n+1}\mathcal{L}_{h}^{n+1}(f^{n+1}) = \varepsilon \underbrace{\left[ \frac{f^{n+1} - f^{(2)}}{\Delta t} - P_X^{n+1}\mathcal{A}_h(f^{(2)}) \right]}_{:= G^{n+1}}.
\end{equation}
In the assumed asymptotic regime, the term $G^{n+1}$ depends only on bounded transport and time-difference operators acting on the smooth solution; thus $\|G^{n+1}\|_{xv}^h = \mathcal{O}(1)$ as $\varepsilon \to 0$.

We now decompose the Fokker--Planck operator image into projected and orthogonal components:
\[
  \big\|\mathcal L_{h}^{n+1}(f^{n+1})\big\|_{xv}^h \le \big\|P_X^{n+1}\mathcal L_{h}^{n+1}(f^{n+1})\big\|_{xv}^h + \big\|(I-P_X^{n+1})\mathcal L_{h}^{n+1}(f^{n+1})\big\|_{xv}^h.
\]
From \eqref{eq:G_derivation}, the projected component is exactly $\varepsilon \|G^{n+1}\|_{xv}^h$. For the orthogonal component, since $f^{n+1}$ lies in the range of $P_X^{n+1}$, Proposition~\ref{prop:fluctuation_bound} bounds the orthogonal component by $\kappa_h \|f^{n+1}\|_{xv}^h$.

To utilize the coercivity estimate, we use the micro--macro decomposition $f^{n+1} = \rho^{n+1} M^{n+1} + g^{n+1}$. We bound the state norm by
\[
  \|f^{n+1}\|_{xv}^h \le \|\rho^{n+1}M^{n+1}\|_{xv}^h + \|g^{n+1}\|_{xv}^h \le M_{\max}^{n+1}\sqrt{L_v}\|\rho^{n+1}\|_x^h + \|g^{n+1}\|_{xv}^h.
\]

We now test the operator against the relative fluctuation $g^{n+1}/M^{n+1}$. Proposition~\ref{prop:coercivity} provides the lower bound $-\langle \mathcal L_{h}^{n+1}(f^{n+1}),\,g^{n+1}/M^{n+1} \rangle_{xv}^h \ge \gamma_h\,\|g^{n+1}/M^{n+1}\|_{xv}^{h,2}$. Conversely, decomposing $\mathcal L_{h}^{n+1}$ into $\varepsilon G^{n+1} + (I-P_X^{n+1})\mathcal L_{h}^{n+1}$ and testing against $g^{n+1}/M^{n+1}$ yields an upper bound:
\begin{align*}
  -\big\langle \mathcal L_{h}^{n+1}(f^{n+1}),\,g^{n+1}/M^{n+1} \big\rangle_{xv}^h
  &\le \big( \varepsilon \|G^{n+1}\|_{xv}^h + \kappa_h \|f^{n+1}\|_{xv}^h \big) \|g^{n+1}/M^{n+1}\|_{xv}^h \\
  &\le \Big( \varepsilon \|G^{n+1}\|_{xv}^h + \kappa_h \big[ M_{\max}^{n+1}\sqrt{L_v}\|\rho^{n+1}\|_x^h + \|g^{n+1}\|_{xv}^h \big] \Big) \|g^{n+1}/M^{n+1}\|_{xv}^h.
\end{align*}
Using the inequality $\|g^{n+1}\|_{xv}^h = \|M^{n+1} (g^{n+1}/M^{n+1})\|_{xv}^h \le M_{\max}^{n+1} \|g^{n+1}/M^{n+1}\|_{xv}^h$, we can absorb the $\|g^{n+1}\|$ term. Assuming $\|g^{n+1}/M^{n+1}\|_{xv}^h > 0$, we divide by $\|g^{n+1}/M^{n+1}\|_{xv}^h$ and rearrange terms:
\[
  (\gamma_h - \kappa_h M_{\max}^{n+1}) \|g^{n+1}/M^{n+1}\|_{xv}^h \le \varepsilon \|G^{n+1}\|_{xv}^h + \kappa_h M_{\max}^{n+1} \sqrt{L_v} \|\rho^{n+1}\|_x^h.
\]
Identifying the amplification factor $\Theta = \gamma_h / (\gamma_h - \kappa_h M_{\max}^{n+1})$, we obtain the bound for $\|g^{n+1}/M^{n+1}\|_{xv}^h$. The distance to equilibrium is simply $\|f^{n+1}-\rho^{n+1}M^{n+1}\|_{xv}^h = \|g^{n+1}\|_{xv}^h \le M_{\max}^{n+1}\|g^{n+1}/M^{n+1}\|_{xv}^h$, which yields \eqref{eq:distance_bound}. The residual bound \eqref{eq:residual_bound} follows by substituting this back into the norm decomposition of $\mathcal{L}_{h}$.
\end{proof}

\end{document}
